\section{Métricas de eficácia e critérios de decisão}
\label{sec:metricas}

As métricas foram escolhidas para acompanhar o funil completo e evitar que apenas vaidade de alcance determine a decisão. Os valores iniciais abaixo são hipóteses de trabalho para o piloto; devem ser recalibrados após a primeira rodada.

\begin{longtable}{p{0.22\textwidth}p{0.29\textwidth}p{0.34\textwidth}}
\toprule
\textbf{Categoria} & \textbf{Métrica} & \textbf{Pergunta de decisão}\\
\midrule
\endhead
Comunicação & Taxa de clique/QR e taxa de cadastro por origem & O canal atrai pessoas do público-alvo para a solução?\\
Comunicação & Custo por cadastro qualificado (CAC inicial) & O alcance é viável para a startup no estágio atual?\\
Comunicação & Compartilhamentos e indicações & A mensagem é compreendida e considerada confiável?\\
\addlinespace
Distribuição & Taxa de conclusão do onboarding & A porta de entrada é simples o suficiente?\\
Distribuição & Tempo até primeiro valor (primeiro lembrete ou agendamento) & O tutor consegue perceber benefício rapidamente?\\
Distribuição & Comparecimento aos atendimentos e uso de créditos & O canal conecta efetivamente o tutor ao cuidado?\\
\addlinespace
Relacionamento & Taxa de leitura e resposta às notificações & Os lembretes são oportunos e úteis?\\
Relacionamento & Retenção em 30 dias e retorno ao serviço & O relacionamento promove continuidade do cuidado?\\
Relacionamento & Satisfação pós-atendimento e taxa de resolução & O suporte gera confiança e resolve a necessidade?\\
\bottomrule
\end{longtable}

\subsection{Regra de aprendizado}
Ao final de cada semana, a equipe deve comparar os canais por conversão e custo, registrar qualitativamente as barreiras encontradas e decidir entre manter, ajustar ou interromper cada experimento. Um canal só deve ser ampliado quando apresentar simultaneamente sinais de adequação ao público, entrega de valor e viabilidade operacional. O aprendizado deve voltar ao desenho do produto, da comunicação e das parcerias.

\section*{Considerações finais}
\addcontentsline{toc}{section}{Considerações finais}
O portfólio proposto combina canais de baixo custo e alta proximidade territorial com uma base digital que organiza a prevenção. A estratégia mantém coerência com a proposta do PetCuida: comunicar com confiança, entregar acesso possível e sustentar uma relação contínua de cuidado. O piloto permitirá substituir hipóteses por evidências antes de ampliar a operação.
